% Generated draft: requires mathematical and source review.
\documentclass[11pt]{article}
\usepackage[margin=25mm]{geometry}
\usepackage{fontspec}
\setmainfont{Latin Modern Roman}
\usepackage{amsmath,amssymb,mathtools,unicode-math}
\setmathfont{Latin Modern Math}
\usepackage{enumitem,needspace,xurl,hyperref}
\setlength{\parindent}{0pt}
\setlength{\parskip}{5pt}
\setlength{\emergencystretch}{3em}
\newcommand{\N}{\mathbb{N}}
\newcommand{\Z}{\mathbb{Z}}
\newcommand{\Q}{\mathbb{Q}}
\newcommand{\R}{\mathbb{R}}
\newcommand{\C}{\mathbb{C}}
\newcommand{\F}{\mathbb{F}}
\newcommand{\A}{\mathbb{A}}
\newcommand{\PP}{\mathbb{P}}
\newcommand{\E}{\mathbb{E}}
\newcommand{\eps}{\varepsilon}
\newcommand{\dd}{\,\mathrm{d}}
\begin{document}
\section*{UnsolvedMath ID: 20002081}
\textbf{Source problem code:} catalogue entry
\subsection{Definitions and mathematical statement}
The source reports Polterovich's suggestion that pullback Hamiltonians from a fibration form a ``maximal torus'' and provide Floer-homology examples. It explicitly leaves the fibration type uncertain and says later locally toric discussion was not recorded. It also has a type error: for \(\Phi:M\to B\), the pullback must use \(h:B\to\R\), not the printed \(H:M\to\R\). Functions are Hamiltonians, not elements of \(\operatorname{Ham}(M)\) until their flows are taken. The quoted torus terminology is therefore not a literal general theorem.

For a precise regular statement, let \((M,\omega)\) be symplectic and \(\Phi:M\to B\) a surjective smooth submersion. Put
\[
 \mathcal A_\Phi=\Phi^*C^\infty(B).
\]
The aid proves that \(\mathcal A_\Phi\) Poisson-commutes if and only if every fiber is coisotropic: for \(V=\ker\dd\Phi\), \(V^\omega\subset V\). If the fibers are connected Lagrangians, the full Poisson centralizer is
\[
 \{F\in C^\infty(M):\{F,h\circ\Phi\}=0\text{ for all }h\in C^\infty(B)\}
       =\mathcal A_\Phi.
\]
Thus this is a maximal Poisson-commutative algebra, not automatically a finite-dimensional torus group. Completeness is additionally needed to turn every Hamiltonian into a globally defined flow; compactness supplies it in the following case.

Let \(M^{2n}\) be compact connected toric symplectic, with effective \(T^n=\R^n/\Z^n\) action and moment map \(\Phi:M\to\Delta\). Choose the period-one action-angle normalization. For smooth \(h\) on the moment polytope, meaning smoothly extendible as required for \(h\circ\Phi\), its interior flow is
\[
 (p,\theta)\longmapsto(p,\theta+t\nabla h(p)).
\]
All such time-one flows form an infinite-dimensional abelian group. The exact kernel of the map \(h\mapsto\varphi_{h\circ\Phi}^1\) consists of
\[
 h(p)=c+\langle m,p\rangle,\qquad c\in\R,\ m\in\Z^n.
\]
The affine-linear subfamily \(h(p)=c+\langle a,p\rangle\), \(a\in\R^n\), alone yields the genuine toric \(T^n\), a maximal Hamiltonian torus. Its lattice depends on the explicitly chosen period normalization.

On the free interior locus, an entire fiber is a one-periodic orbit family precisely at
\[
 \nabla h(p)\in\Z^n.
\]
The zero integer vector gives constant trajectories. If such a resonant \(p\) is isolated and \(\operatorname{Hess}h(p)\) is nonsingular, its fiber is Morse--Bott: the kernel of the linearized time-one map minus the identity is exactly the tangent space of that torus. Degenerate Hessians can require a larger critical-family analysis. Boundary faces have isotropy and are outside this interior criterion. These statements clarify proposed Floer examples but do not compute Floer homology or supply the missing singular-fibration hypotheses.
\subsection{Short English statement}
The corrected pullback algebra commutes precisely for coisotropic fibers and is maximal as a Poisson algebra for connected Lagrangian fibers. In a compact toric example its general flow group is infinite-dimensional; only affine-linear Hamiltonians give the ordinary torus. Integral-gradient resonances describe interior periodic tori for Floer analysis.
\subsection{Sources}
\begin{itemize}
\item Local UnsolvedMath registry, permanent ID \(20002081\), source code \texttt{catalogue entry}; canonical collection \texttt{aim-geometry-notes.json}, notes[418]. Statement status \texttt{reconstructed\_unverified}; clean statement null.
\item Michael Hutchings, with help from Yasha Eliashberg and John Etnyre, ``Holomorphic curves in contact geometry,'' AIM workshop problem list, version 15 October 2003, Question \(1.23\): \url{https://aimath.org/WWN/contactgeom2/contactgeom2.pdf}. Polterovich is the suggestion's named source. The supplied verification says PDF and HTML share the type error; neither was opened in this session.
\item Supplied research attempt \(1\), result \texttt{context\_only}, run status \texttt{partially\_solved}, rigor label ``proved,'' novelty confidence ``low.'' Registry status ``open,'' classification \texttt{OPEN-TRIAGE}, suggested difficulty \texttt{L2}, updated 14 August 2026. Original remarks and literature are empty; no full authoritative theorem references or proposer/year registry values are supplied.
\end{itemize}
\end{document}
